\scanpage{22}
Cette section mérite le nom de \uline{invariant de
Hasse} de $A$. L'ens des zéros de cette
section est \unsure{exactmt|exactement} l'ens des $s$ tels que
$A_s$ soit \uline{non} \uline{ordinaire}.

Quelles sont les relations entre $\omega_A$ et $\omega_{A^*}$ ? On
a, grâce à la suite exacte
\[
\unsure{0|\cdot}\mathbin{\unsure{-|\to}}\underline t_A^\vee\to\mathrm{DR}^1(A/S)\mathbin{-}\underline t_{A^*}\mathbin{-}0
\]
la relation
\diagram{\begin{tikzcd}[ampersand replacement=\&,column sep=small,row sep=tiny]
\det t_A^\vee\otimes\det\underline t_{A^*}\arrow[r,phantom,"\simeq"] \& \det\mathrm{DR}^1(A/S)\arrow[d,phantom,"\unsure{\simeq|=|\sim}"{description,rotate=90}]\\
\& \mathrm{DR}^{2g}(A/S)\simeq\underline O_S
\end{tikzcd}}
donc on aura
\[
\omega_A\cdot\omega_{A^*}^\vee\simeq\underline O_S
\]
i.e. on a un isom can.
\[
\omega_A\simeq\omega_{A^*}
\]
Je dis que \uline{cet isom est compatible avec les
p-structures}, \deleted{\unsure{et|---}} \uline{donc transforme} $\sigma_A$ \uline{en} $\sigma_{A^*}$, \uline{et
dans le cas} $A$ \uline{ordinaire}, induit un isom
\[
\operatorname{d\acute et}({}_pA)_{\mathrm{\acute et}}\simeq\operatorname{d\acute et}({}_p A^*)_{\mathrm{\acute et}} .
\]
% archive sheet number: 21
